\emph{Predictive bounds for the prior-centre comparison.}\label{app:r3:centres}
For $\gamma,v>0$, write $A_\gamma=F+\gamma I$. The zero-mean-prior shift splits as
\[
 \delta_\gamma=\delta_{L,\gamma}+\delta_{P,\gamma},\qquad
 \delta_{L,\gamma}=-A_\gamma^{-1}J^\top r,\qquad
 \delta_{P,\gamma}=-\gamma A_\gamma^{-1}\theta.
\]
In the factorisation above, their coordinates are
$-s_i(V^\top r)_i/(s_i^2+\gamma)$ and
$-\gamma\theta_{\rm modes,i}/(s_i^2+\gamma)$; the QR-complement contributions
are $0$ and $-Q_cQ_c^\top\theta$, respectively. The total norm includes their
cross term. Replacing the prior by $\vartheta\sim\mathcal N(\theta,\alpha^{-1}I)$
removes $\delta_{P,\gamma}$ and leaves the covariance unchanged.

Define $\mathcal K_\gamma=r^\top J A_\gamma^{-1}J^\top r$ and
$c_\gamma=\sqrt{\mathcal K_\gamma/v}$. Since
$\delta_{L,\gamma}^\top A_\gamma\delta_{L,\gamma}=\mathcal K_\gamma$,
Cauchy--Schwarz in the $A_\gamma$ metric gives the pointwise bound
\begin{equation}\label{eq:r3:centre-bound}
 |m^\theta_\gamma(x)-u_\theta(x)|
 =|g(x)^\top\delta_{L,\gamma}|
 \le\sigma(x)c_\gamma,\qquad
 \mathcal K_\gamma=\sum_i\frac{s_i^2(V^\top r)_i^2}{s_i^2+\gamma}.
\end{equation}
Let $E_r$ and $\Sigma_r$ denote trained error RMS and uncertainty RMS, and let
$q_{r,i}$ be the regional mode energies in \eqref{eq:r2-regional-modal}.
The regional RMS triangle inequality also gives
\begin{equation}\label{eq:r3:regional-bound}
 B_r=\min\!\left\{c_\gamma\Sigma_r,
              \sum_i|\delta_{L,\gamma,i}|\sqrt{q_{r,i}}\right\},\qquad
 \max(0,E_r-B_r)\le E^\theta_r\le E_r+B_r.
\end{equation}
For $E^\theta_{\mathcal B}>0$ and $\Rrms>0$,
$\rho^\theta\ge\max(0,E_{\mathcal S}-B_{\mathcal S})/
[(E_{\mathcal B}+B_{\mathcal B})\Rrms]$.
The prior-only prediction differs from
the archived zero-mean affine mean by $g^\top\delta_L$; its full-grid error is
therefore at least $E^{\rm zero}_{\rm full}-B_{\rm full}$. This attributes
worsening through a prediction-level bound rather than parameter norms.
If $C_c(t)$ denotes empirical coverage at $t$ standard deviations, then
\[
 C_{\rm trained}(t-c_\gamma)\le C_\theta(t)
       \le C_{\rm trained}(t+c_\gamma),\qquad t\ge c_\gamma,
\]
by the pointwise triangle inequality.

% Derived from round3_local/outputs/summary.json, 12 October 2026.
The saved-array calculation reproduces all $401$ total displacement norms to
$4.44\times10^{-16}$ relative and regional variances to $6.73\times10^{-14}$.
Across those settings, $\|\delta_\gamma\|/\|\theta\|$ ranges from approximately
$0.9889$ to $31.17$; these displacements do not validate a nonlinear first-order approximation.
At the baseline, $\|\delta\|/\|\theta\|$ is $[0.997775,0.999576]$, whereas
$\|\delta_L\|/\|\theta\|$ is $[8.73\times10^{-6},1.5413\times10^{-4}]$;
the QR-complement norm alone is $[0.976675,0.985087]\|\theta\|$.
Using \eqref{eq:r3:centre-bound} and \eqref{eq:r3:regional-bound} gives
$c_\gamma\in[0.007489,0.090775]$ and, after the
regional refinement, establishes a $\theta$-centred allocation deficit in $16/17$
runs from bounds alone. Its seed-$4$ lower bound $0.4878$ is inconclusive;
direct predictions establish the deficit in all $17$ (Section~\ref{sec:ceiling}).
The prior-only prediction is worse in all $14$ baseline cases where the
zero-mean affine prediction worsens.
The $\theta$-centred one-sigma coverage lower bounds are $0.7925$ in sharp and
$0.997196$ in bulk, so both regions overcover in all $17$ runs. Bulk $1.96$-sigma
coverage is $1$; sharp $1.96$-sigma coverage is above nominal in $12$ runs and
below nominal in the other $5$. Their upper bounds are $0.924$, $0.883$,
$0.90875$, $0.91225$ and $0.896$ for Burgers--MLP seeds $1,2,3$ and
Burgers--Fourier seeds $0,6$, respectively. At the conditional MLE,
$c_\gamma\in[75.428,79.910]$ is uninformative; Table~\ref{tab:r2-centres} instead
reports the directly evaluated alternative centre.


\emph{Coverage envelopes at refined contrast-matching roots.}\label{app:r3:root-envelope}
At unit noise, each pointwise variance has the positive modal representation
$q_\gamma(x)=\sum_i a_i(x)/(\lambda_i+\gamma)$, including the QR complement
as a zero-eigenvalue mode, with $a_i(x),\lambda_i\ge0$.
For positive endpoint variances at saved precisions
$0<l\le\gamma\le h$, $l<h$, and
$t=(\gamma-l)/(h-l)$, positivity gives
\begin{equation}\label{eq:r3:variance-envelope}
 L_\gamma=\frac{(1-t)lq_l+thq_h}{\gamma}
 \le q_\gamma\le
 U_\gamma=\left(\frac{1-t}{q_l}+\frac{t}{q_h}\right)^{-1}.
\end{equation}
For the lower bound in \eqref{eq:r3:variance-envelope}, $\gamma q_\gamma$ is concave because every
$\gamma/(\lambda_i+\gamma)$ is concave. For the upper bound,
$q'=-\sum_i a_i/(\lambda_i+\gamma)^2$ and
$q''=2\sum_i a_i/(\lambda_i+\gamma)^3$ satisfy
$(q')^2\le qq''/2$ by Cauchy--Schwarz. Hence
$(1/q)''=[2(q')^2-qq'']/q^3\le0$, and taking the reciprocal of its chord gives
$U_\gamma$. Both envelopes decrease with $\gamma$ since $q_l\ge q_h$.

For a numerically refined root bracket $[a,b]\subset[l,h]$, use $L_b$ and
$U_a$ to retain root-location uncertainty. The higher empirical bulk-$0.68$
quantile $Q_{\mathcal B}$ is monotone, so
\[
 k_{\min}=Q_{\mathcal B}(|e|/\sqrt{U_a}),\qquad
 k_{\max}=Q_{\mathcal B}(|e|/\sqrt{L_b})
\]
enclose the multiplier $k_\gamma$ that calibrates bulk coverage at the root.
For sharp coverage at threshold $z\ge0$, this yields
\begin{equation}\label{eq:r3:coverage-envelope}
 \left\langle\mathbf1\{|e|\le z k_{\min}\sqrt{L_b}\}\right\rangle_{\mathcal S}
 \le C_{\mathcal S,\gamma}(z)\le
 \left\langle\mathbf1\{|e|\le z k_{\max}\sqrt{U_a}\}\right\rangle_{\mathcal S}.
\end{equation}
These analytic bounds use the returned spectral model and saved endpoint fields;
they are not directed-rounding certificates for the original Jacobian.
Direct pointwise root evaluations confirm the enclosures in Table~\ref{tab:r3-roots}.
Root refinement concerns the five observed
trained-centre brackets, without asserting uniqueness or global optimality.
Appendix~\ref{oa:lean} maps the Lean verification of these analytic bounds.
The evaluated run counts and refined brackets remain numerical results.

% Generated by scripts/make_paper_tables.py; old bounds and new direct evaluations.
\begin{table}[t]
\centering\small\setlength{\tabcolsep}{4pt}
\begin{tabular}{@{}rrrcr c@{}}
\toprule
 & & \multicolumn{2}{c}{Sharp $1\sigma$ coverage} & \multicolumn{2}{c}{Sharp $1.96\sigma$ coverage} \\
\cmidrule(lr){3-4}\cmidrule(l){5-6}
seed & refined $\gamma$ & direct & envelope & direct & envelope \\
\midrule
$0$ & $0.02105515$ & $0.3275$ & $[0.323,0.333]$ & $0.78275$ & $[0.7815,0.78475]$ \\
$4$ & $0.73626974$ & $0.427$ & $[0.41975,0.43375]$ & $0.852$ & $[0.84525,0.862]$ \\
$5$ & $2.0476385\times10^{-5}$ & $0.30325$ & $[0.29775,0.315]$ & $0.67725$ & $[0.6755,0.67925]$ \\
$6$ & $0.00163541$ & $0.52325$ & $[0.5225,0.526]$ & $0.847$ & $[0.8465,0.84775]$ \\
$7$ & $3.0590419\times10^{-5}$ & $0.55725$ & $[0.554,0.563]$ & $0.8235$ & $[0.82125,0.82675]$ \\
\bottomrule
\end{tabular}
\caption{Direct sharp coverage and saved envelopes at five trained Burgers--MLP
roots. Displayed precisions are rounded; root brackets have relative width below
$5.3\times10^{-13}$. Bulk calibration uses the higher empirical $0.68$ quantile.
Envelopes include root and calibration uncertainty through
\eqref{eq:r3:coverage-envelope}; all direct values lie inside.}
\label{tab:r3-roots}
\end{table}

